\thispagestyle{empty}

\noindent{\sffamily\bfseries\large\color{Ink} How to use this book}
\grule
\vspace{6pt}

Every topic gets exactly one page. Nothing runs over, so you can open the book
flat at a single table and drill it.

\tcap{The colour code}

Gender keeps the same colour on every page. Learn the colour with the word and
the tables start to read themselves.

\begin{gendertbl}{}
  Gender  & maskulin & feminin & neutrum & Plural \\
  Article & der      & die     & das     & die    \\
  Example & der Mann & die Frau & das Kind & die Leute \\
\end{gendertbl}

\tcap{The four cases}

\begin{gtbl}{colspec={X[0.9,l] X[1.5,l] X[2,l]}}
  Case & Marks & Rough test \\
  Nominativ & the subject & \emph{who or what} is doing it \\
  Akkusativ & the direct object & \emph{whom or what} is affected \\
  Dativ & the indirect object & \emph{to or for whom} \\
  Genitiv & possession & \emph{whose} \\
\end{gtbl}

\begin{gnote}
Only the \textbf{masculine singular} changes between nominative and accusative
(\textit{der}~$\rightarrow$~\textit{den}). Feminine, neuter and plural look
identical in those two cases. Half of German case anxiety disappears once that
sinks in.
\end{gnote}

\tcap{Symbols}

\begin{flattbl}{colspec={X[0.7,l] X[2.3,l]}}
  \en{-en}    & a violet ending is the bit that changes \\
  \xx         & no such form exists \\
  \eng{gloss} & grey italics are an English gloss or a hint \\
\end{flattbl}
